% Native hgraph authoring; original AimPaper statements are checked verbatim against AimPaper/main.tex.
\section{Selected classes, survival ranges and incoming differentials}

% AimPaper: 7.6 fact 2156 2167
\begin{fact}[AimPaper Fact 7.6]
\label{paper:7.6}
\source{AimPaper Fact 7.6}
\level{coarse}
From \cite{LWX126}, Fact 7.6.\par
% BEGIN AIMPAPER STATEMENT
\label{fact:theta5sqAF}
\begin{enumerate}
   % \item[]
    \item $x_{126,8,4} + x_{126,8}$ survives to the $E_6$-page.
    
    \item $h_1h_4x_{109,12}$ is a permanent cycle, and can only be killed by 
    $$d_6(x_{126,8,4} + x_{126,8}) \ \text{or} \ d_{12}(h_6^2).$$
    
    \item $h_0^2x_{124,8}$ survives to the $E_\infty$-page.

    \item In $\Ext_A^{25,125+25}$, the element $g^4\Delta h_1g$ is the only one that survives to the classical $E_5$-page.
\end{enumerate}
% END AIMPAPER STATEMENT
\uses{sem:Computation:table-bases,sem:Computation:table-certification}
% Blueprint note: Mathematical dependency review: [AimPaper source, lines 2142-2177](paper/main.tex.html#L2142). The four survival and incoming-differential facts are read from the stated sphere tables and the near-126 chart, rather than deduced from the permanence theorem.



% Correspondence: public counterparts; scope is specified below.
% Blueprint note: The public statements separate the four computational inputs: survival of $W$ to a finite page, the cycle condition and possible incoming differentials for $T$, survival of $U$, and the restriction in high filtration.
% Blueprint note: \par This entry is bound to the existing Lean declarations. Its status reports their local completion; the scope of correspondence with the article is specified above.
% Blueprint note: Formal counterparts: \texttt{KIP126.Computation.Standard.w\_to\_e6}; \texttt{KIP126.Computation.Standard.t\_permanent\_cycle}; \texttt{KIP126.Computation.Standard.t\_only\_incoming}; \texttt{KIP126.Computation.Standard.u\_permanent}; \texttt{KIP126.Computation.Standard.high\_e5}
\lean{KIP126.Computation.Standard.w_to_e6,KIP126.Computation.Standard.t_permanent_cycle,KIP126.Computation.Standard.t_only_incoming,KIP126.Computation.Standard.u_permanent,KIP126.Computation.Standard.high_e5}

% Blueprint note: Original article number 7.6. \href{paper/main.tex.html#L2156}{Original statement, hypotheses, and context} · \href{paper/2412.10879.pdf}{Article PDF}
\end{fact}
\begin{proof}
\textit{(See \cite{LWX126}, Fact 7.6.)} \href{paper/main.tex.html#L2156}{Original statement and context}.

Read the sphere computations in \cite{LWX126}, Figure 1 and the Appendix tables in stems 124, 125 and 126, using the stated bidegrees. The row for $x_{126,8,4}+x_{126,8}$ has no outgoing differential before $d_6$. The row for $h_1h_4x_{109,12}$ has no outgoing differential, and the source rows leave only the displayed $d_6$ and $d_{12}$ as incoming possibilities. In particular the stem-126 filtration-six source supports its nonzero $d_3$ and is not an additional incoming candidate.

The stem-124 row gives the infinite-cycle condition for $h_0^2x_{124,8}$. The stem-125 filtration-25 rows leave only $g^4\Delta h_1g$ on $E_5$. These are certified table inputs, not deductions from the new permanence theorem; their fixed-page and fixed-label verification is the computational responsibility retained in this chapter.
\end{proof}

% AimPaper: 7.13 fact 2377 2385
\begin{fact}[AimPaper Fact 7.13]
\label{paper:7.13}
\source{AimPaper Fact 7.13}
\level{coarse}
From \cite{LWX126}, Fact 7.13.\par
% BEGIN AIMPAPER STATEMENT
\label{fact:x_123_9}\label{fact:x1239}
    \begin{enumerate}

   \item  $x_{123,9}+h_0x_{123,8}$ survives to the Adams $E_{12}$-page, and is not killed by any classical differential.
      % \item In Ext, $h_1 \cdot (x_{123,9}+h_0x_{123,8}) = h_1 x_{123,9}$.
    \item We have a classical nonzero differential 
    $$d_2 (x_{125,8}) = h_1 (x_{123,9}+h_0x_{123,8}) + h_0^2x_{124,8}.$$
   \end{enumerate}
% END AIMPAPER STATEMENT
\uses{sem:Computation:table-bases,sem:Computation:table-certification}
% Blueprint note: Mathematical dependency review: [AimPaper source, lines 2375-2385](paper/main.tex.html#L2375). The survival range and specified d2 differential are independent inputs from the stem-123 and stem-125 sphere tables.



% Correspondence: public counterparts; scope is specified below.
% Blueprint note: The public statements give survival of $V$ to $E_{12}$, the absence of incoming differentials, and the specified nonzero $d_2$ differential.
% Blueprint note: \par This entry is bound to the existing Lean declarations. Its status reports their local completion; the scope of correspondence with the article is specified above.
% Blueprint note: Formal counterparts: \texttt{KIP126.Computation.Standard.v\_to\_e12}; \texttt{KIP126.Computation.Standard.v\_not\_hit}; \texttt{KIP126.Computation.Standard.d2\_x125\_8}
\lean{KIP126.Computation.Standard.v_to_e12,KIP126.Computation.Standard.v_not_hit,KIP126.Computation.Standard.d2_x125_8}

% Blueprint note: Original article number 7.13. \href{paper/main.tex.html#L2377}{Original statement, hypotheses, and context} · \href{paper/2412.10879.pdf}{Article PDF}
\end{fact}
\begin{proof}
\textit{(See \cite{LWX126}, Fact 7.13.)} \href{paper/main.tex.html#L2377}{Original statement and context}.

Use the stem-123 and stem-125 tables in the Appendix of \cite{LWX126}. The former records that $x_{123,9}+h_0x_{123,8}$ supports no differential before $d_{12}$ and has no incoming differential. In the latter, the row for $x_{125,8}$ records the nonzero $d_2$ with value $h_1(x_{123,9}+h_0x_{123,8})+h_0^2x_{124,8}$. The target is read in its $E_2$ basis, so both summands are retained. Transport of these table entries to the selected sphere spectral sequence remains the certified computation input.
\end{proof}

% AimPaper: 7.15 fact 2466 2468
\begin{fact}[AimPaper Fact 7.15]
\label{paper:7.15}
\source{AimPaper Fact 7.15}
\level{coarse}
From \cite{LWX126}, Fact 7.15.\par
% BEGIN AIMPAPER STATEMENT
\label{fact:h_0_sq_mul_x_125_9_2}\label{fact:h02x1259}
    The element $h_0^2x_{125,9,2}$ survives to the Adams $E_5$-page, and is not killed by any classical differential.
% END AIMPAPER STATEMENT
\uses{sem:Computation:table-bases,sem:Computation:table-certification}
% Blueprint note: Mathematical dependency review: [AimPaper source, lines 2464-2468](paper/main.tex.html#L2464). The stated E5 survival and absence of incoming hits are read from the stem-125 sphere table.



% Correspondence: public counterparts; scope is specified below.
% Blueprint note: The public statements give survival of $Y$ to $E_5$ and exclude incoming differentials.
% Blueprint note: \par This entry is bound to the existing Lean declarations. Its status reports their local completion; the scope of correspondence with the article is specified above.
% Blueprint note: Formal counterparts: \texttt{KIP126.Computation.Standard.y\_to\_e5}; \texttt{KIP126.Computation.Standard.y\_not\_hit}
\lean{KIP126.Computation.Standard.y_to_e5,KIP126.Computation.Standard.y_not_hit}

% Blueprint note: Original article number 7.15. \href{paper/main.tex.html#L2466}{Original statement, hypotheses, and context} · \href{paper/2412.10879.pdf}{Article PDF}
\end{fact}
\begin{proof}
\textit{(See \cite{LWX126}, Fact 7.15.)} \href{paper/main.tex.html#L2466}{Original statement and context}.

The stem-125 table in the Appendix of \cite{LWX126}, at the row $h_0^2x_{125,9,2}$, excludes incoming differentials and excludes outgoing differentials before $d_5$. Hence this label reaches $E_5$ and is not killed by a classical incoming differential. No claim about the value of its possible $d_5$ follows from this table entry; that value is deliberately not assumed here.
\end{proof}

% AimPaper: 7.19 fact 2613 2615
\begin{fact}[AimPaper Fact 7.19]
\label{paper:7.19}
\source{AimPaper Fact 7.19}
\level{coarse}
From \cite{LWX126}, Fact 7.19.\par
% BEGIN AIMPAPER STATEMENT
\label{fact:h_1_mul_x_121_7}\label{fact:h1x1217}
    The element $h_1x_{121,7}$ survives to the Adams $E_6$-page, and is not killed by any classical differential.
% END AIMPAPER STATEMENT
\uses{sem:Computation:table-bases,sem:Computation:table-certification}
% Blueprint note: Mathematical dependency review: [AimPaper source, lines 2611-2615](paper/main.tex.html#L2611). The E6 survival and absence of incoming differentials come from the stem-122 sphere table.



% Correspondence: public counterparts; scope is specified below.
% Blueprint note: The public statements give survival of $X$ to $E_6$ and exclude incoming differentials.
% Blueprint note: \par This entry is bound to the existing Lean declarations. Its status reports their local completion; the scope of correspondence with the article is specified above.
% Blueprint note: Formal counterparts: \texttt{KIP126.Computation.Standard.x\_to\_e6}; \texttt{KIP126.Computation.Standard.x\_not\_hit}
\lean{KIP126.Computation.Standard.x_to_e6,KIP126.Computation.Standard.x_not_hit}

% Blueprint note: Original article number 7.19. \href{paper/main.tex.html#L2613}{Original statement, hypotheses, and context} · \href{paper/2412.10879.pdf}{Article PDF}
\end{fact}
\begin{proof}
\textit{(See \cite{LWX126}, Fact 7.19.)} \href{paper/main.tex.html#L2613}{Original statement and context}.

In the stem-122 table in the Appendix of \cite{LWX126}, the row for $h_1x_{121,7}$ has no outgoing differential before $d_6$. The same computation excludes incoming differentials to this class. These two row assertions are respectively survival to $E_6$ and the absence of classical incoming hits, with the same standard label. No inference from the later $\nu$-extension is needed.
\end{proof}

\section{Permanent cycles in stem 122}

% AimPaper: 7.21 fact 2689 2696
\begin{fact}[AimPaper Fact 7.21]
\label{paper:7.21}
\source{AimPaper Fact 7.21}
\level{coarse}
From \cite{LWX126}, Fact 7.21.\par
% BEGIN AIMPAPER STATEMENT
\label{fact:stem122}
\begin{enumerate}
    \item[] 
    \item The element $h_6 M d_0$ survives to the Adams $E_\infty$-page.
    \item The element $h_5x_{91,11}$ survives to the Adams $E_\infty$-page.
\end{enumerate}
% END AIMPAPER STATEMENT
\uses{sem:Computation:table-bases,sem:Computation:table-certification}
% Blueprint note: Mathematical dependency review: [AimPaper source, lines 2683-2696](paper/main.tex.html#L2683). Both permanent-cycle inputs are extracted from the stem-122 table before the proof of Proposition 7.9.



% Correspondence: public counterparts; scope is specified below.
% Blueprint note: The public statements give nonzero permanent survival of $P=h_6Md_0$ and $Q=h_5x_{91,11}$.
% Blueprint note: \par This entry is bound to the existing Lean declarations. Its status reports their local completion; the scope of correspondence with the article is specified above.
% Blueprint note: Formal counterparts: \texttt{KIP126.Computation.Standard.p\_permanent}; \texttt{KIP126.Computation.Standard.q\_permanent}
\lean{KIP126.Computation.Standard.p_permanent,KIP126.Computation.Standard.q_permanent}

% Blueprint note: Original article number 7.21. \href{paper/main.tex.html#L2689}{Original statement, hypotheses, and context} · \href{paper/2412.10879.pdf}{Article PDF}
\end{fact}
\begin{proof}
\textit{(See \cite{LWX126}, Fact 7.21.)} \href{paper/main.tex.html#L2689}{Original statement and context}.

The stem-122 table in the Appendix of \cite{LWX126} records no outgoing differential on either $h_6Md_0$ or $h_5x_{91,11}$, giving the two infinite-cycle assertions. The corresponding certified sphere and product computations are the inputs used when these labels are lifted in the contradiction argument. This fact states the two permanence computations; it does not by itself establish a general rule for lifting relations from a finite $\lambda$ quotient.
\end{proof}
